% ============================================================================
% docs/notes.tex -- Living project notebook.
%
% MAINTENANCE RULES (see also CLAUDE.md, section "Project notebook"):
%   - This file is updated incrementally as the project progresses, in the
%     same change as the code it documents. It is not a one-off writeup.
%   - Every module implemented in riskparity/ should get its formulas,
%     reasoning, and a short annotated code excerpt added to the matching
%     subsection of Section "Implementation Log" -- do not leave a module
%     documented only in code comments.
%   - Figures are pulled from results/figures/ via \IfFileExists guards, so
%     the document always compiles even before the pipeline has produced
%     any output; once a figure exists on disk it appears automatically on
%     the next compile.
%   - Keep derivations self-contained: a reader should be able to follow the
%     math here without cross-referencing the paper unless it's explicitly
%     cited for a result that isn't re-derived.
%   - Compile with: latexmk -pdf docs/notes.tex   (or pdflatex, run twice
%     for the table of contents / references to resolve).
% ============================================================================

\documentclass[11pt]{article}

\usepackage[a4paper,margin=2.5cm]{geometry}
\usepackage{amsmath,amssymb,amsthm}
\usepackage{graphicx}
\usepackage{booktabs}
\usepackage[dvipsnames]{xcolor}
\usepackage{listings}
\usepackage{hyperref}
\usepackage{enumitem}
\usepackage{todonotes}

\hypersetup{colorlinks=true, linkcolor=blue!50!black, citecolor=blue!50!black, urlcolor=blue!50!black}

\newtheorem{theorem}{Theorem}
\newtheorem{definition}{Definition}
\newtheorem{proposition}{Proposition}

\lstdefinestyle{python}{
    language=Python,
    basicstyle=\ttfamily\small,
    keywordstyle=\color{blue!60!black}\bfseries,
    commentstyle=\color{gray},
    stringstyle=\color{OliveGreen},
    numbers=left,
    numberstyle=\tiny\color{gray},
    showstringspaces=false,
    breaklines=true,
    frame=single,
    tabsize=4,
}
\lstset{style=python}

\newcommand{\Var}{\mathrm{VaR}}
\newcommand{\CVar}{\mathrm{CVaR}}

\title{Risk Parity / ERC vs.\ Component CVaR Parity\\[4pt]\large Project Notebook: Finance Concepts, Mathematics, and Implementation}
\author{Quant Risk Projects}
\date{Started \today\ --- updated continuously alongside \texttt{riskparity/}}

\begin{document}
\maketitle

\begin{abstract}
This notebook is the running mathematical and financial companion to the
\texttt{risk-parity-cvar} codebase. It collects, in one place and in
increasing detail as the project advances: the finance concepts involved
(risk parity, tail risk, diversification), the mathematics behind each
risk measure and solver (with derivations, not just citations), the
corresponding Python implementation with commentary, and the empirical
results once they exist. See \texttt{README.md} for the project pitch and
\texttt{PROJECT\_BRIEF.md} for scope/design rationale; this document is
where the \emph{content} -- the actual finance and math -- lives.
\end{abstract}

\tableofcontents

% ============================================================================
\section{Motivation}
% ============================================================================

Volatility-based risk parity treats risk as symmetric and approximately
Gaussian: an asset's contribution to portfolio risk depends only on its
variance and covariances with other assets. In practice, tail risk is not
proportional to variance -- some assets have fatter left tails than their
volatility alone suggests (e.g.\ equities vs.\ long-duration government
bonds). The central question of this project:

\begin{quote}
\itshape Does equalizing \textbf{Component CVaR} contributions instead of
volatility contributions produce a materially different, more
tail-robust portfolio -- and if so, when does it matter most?
\end{quote}

Answering this requires two parity engines built on the same mathematical
principle (Euler risk decomposition, Section~\ref{sec:euler}) but applied
to two different risk measures: one with a closed form (volatility,
Section~\ref{sec:vol-erc}) and one without (CVaR,
Section~\ref{sec:cvar-parity}).

% ============================================================================
\section{Mathematical Foundations: Euler Risk Decomposition}
\label{sec:euler}
% ============================================================================

\subsection{Positive homogeneity}

\begin{definition}[Positive homogeneity of degree 1]
A function $\rho : \mathbb{R}^n \to \mathbb{R}$ is \textbf{positively
homogeneous of degree 1} if
\[
\rho(t\, w) = t \, \rho(w) \qquad \text{for all } t > 0,\ w \in \mathbb{R}^n .
\]
\end{definition}

Portfolio volatility and portfolio CVaR both have this property: if every
position in the portfolio is scaled by $t>0$ (before renormalizing weights
to sum to 1, or simply thinking of $w$ as dollar exposures), the risk
measure scales by exactly $t$. This is the property that makes an
\emph{exact} risk decomposition possible.

\begin{theorem}[Euler's homogeneous function theorem]
If $\rho$ is positively homogeneous of degree 1 and differentiable, then
\[
\rho(w) = \sum_{i=1}^n w_i \, \frac{\partial \rho}{\partial w_i}(w).
\]
\end{theorem}

\begin{proof}[Proof sketch]
Differentiate both sides of $\rho(tw) = t\rho(w)$ with respect to $t$:
\[
\sum_i w_i \, \frac{\partial \rho}{\partial w_i}(tw) = \rho(w).
\]
Setting $t = 1$ gives the result.
\end{proof}

\subsection{Risk contributions}

\begin{definition}[Marginal, absolute, and percentage risk contribution]
\label{def:rc}
For a risk measure $\rho(w)$ satisfying the conditions above:
\[
\text{MRC}_i = \frac{\partial \rho}{\partial w_i}(w), \qquad
\text{RC}_i = w_i \cdot \text{MRC}_i, \qquad
\text{\%RC}_i = \frac{\text{RC}_i}{\rho(w)} .
\]
By Euler's theorem, $\sum_i \text{RC}_i = \rho(w)$ exactly, and
$\sum_i \text{\%RC}_i = 1$. This holds for \emph{any} weights $w$, not
just at an optimum -- it is an identity, not a property of a particular
portfolio.
\end{definition}

\begin{definition}[Risk Parity / Equal Risk Contribution (ERC)]
Given risk budgets $b_i \geq 0$ with $\sum_i b_i = 1$ (defaulting to
$b_i = 1/n$), a portfolio $w$ is \textbf{risk-parity} (or budget-balanced)
with respect to $\rho$ if
\[
\text{RC}_i(w) = b_i \cdot \rho(w) \qquad \text{for all } i.
\]
\end{definition}

This identity and definition are the backbone of the entire project: both
the volatility engine and the CVaR engine implement a solver that searches
for weights satisfying this condition, for two different choices of
$\rho$. It is verified directly by
\texttt{tests/test\_contributions.py} (see
Section~\ref{sec:vol-erc-impl}) -- if a code change ever breaks
$\sum_i \text{RC}_i = \rho(w)$, that is a correctness bug, not something to
paper over with a looser tolerance.

% ============================================================================
\section{Volatility-Based Risk Parity (Vol-ERC)}
\label{sec:vol-erc}
% ============================================================================

\subsection{Portfolio volatility and its closed-form gradient}

Let $\Sigma$ be the (annualized) covariance matrix of asset returns and
$w$ the portfolio weights. Portfolio volatility is
\[
\sigma_p(w) = \sqrt{w^\top \Sigma w},
\]
which is positively homogeneous of degree 1 (scaling $w$ by $t$ scales
$\sigma_p$ by $t$, since the quadratic form scales by $t^2$ and the square
root brings it back to $t$). Its gradient is closed-form:
\[
\frac{\partial \sigma_p}{\partial w_i} = \frac{(\Sigma w)_i}{\sigma_p(w)}
\quad\Longrightarrow\quad
\text{MRC}_i = \frac{(\Sigma w)_i}{\sigma_p(w)}, \qquad
\text{RC}_i = \frac{w_i (\Sigma w)_i}{\sigma_p(w)} .
\]
This is the reason volatility-ERC is analytically tractable: no
simulation or scenario set is required, only $\Sigma$.

\subsection{The ERC optimization problem}

The direct formulation --- minimize the dispersion of risk contributions
---
\[
\min_{w \geq 0,\ \mathbf{1}^\top w = 1} \ \sum_{i,j} \big(\text{RC}_i(w) - \text{RC}_j(w)\big)^2
\]
is not convex in general. Maillard, Roncalli \& Teïletche (2010) show
that for equal budgets $b_i = 1/n$, the ERC solution can instead be
obtained from the convex proxy problem
\[
\min_{w \geq 0} \ \tfrac{1}{2} w^\top \Sigma w - \sum_i b_i \ln w_i,
\]
which is a log-barrier-regularized quadratic form: convex (the Hessian of
$-\sum_i b_i \ln w_i$ is positive definite on the positive orthant, and
$\tfrac12 w^\top \Sigma w$ is convex since $\Sigma \succeq 0$), hence has
a unique minimizer up to scale. Solving it and then renormalizing
$w \leftarrow w / \mathbf{1}^\top w$ yields the ERC portfolio. This convex
reformulation is what makes volatility-ERC the ``easy'' side of the
project, in contrast to CVaR-parity (Section~\ref{sec:cvar-parity}).

\subsection{Two-asset closed form (sanity check case)}
\label{sec:two-asset}

For two uncorrelated assets with volatilities $\sigma_1, \sigma_2$, the
ERC solution has a simple closed form:
\[
w_i \propto \frac{1}{\sigma_i} \quad\Longrightarrow\quad
w_1 = \frac{1/\sigma_1}{1/\sigma_1 + 1/\sigma_2}, \qquad
w_2 = \frac{1/\sigma_2}{1/\sigma_1 + 1/\sigma_2}.
\]
This is the case used in
\texttt{tests/test\_contributions.py::test\_volatility\_contributions\_two\_uncorrelated\_assets\_known\_case}
and in \texttt{tests/test\_erc.py} to validate the general solver against
a known answer.

\subsection{Implementation}
\label{sec:vol-erc-impl}

\texttt{riskparity/risk/contributions.py::volatility\_risk\_contributions}
is the first (and, so far, only) fully implemented function in the
codebase, and serves as the reference pattern for everything else:

\begin{lstlisting}
def volatility_risk_contributions(weights: np.ndarray, cov: np.ndarray) -> RiskContributions:
    weights = np.asarray(weights, dtype=float)
    cov = np.asarray(cov, dtype=float)

    portfolio_variance = weights @ cov @ weights
    sigma_p = np.sqrt(portfolio_variance)

    marginal = (cov @ weights) / sigma_p
    contribution = weights * marginal
    pct_contribution = contribution / sigma_p

    return RiskContributions(
        total_risk=float(sigma_p),
        marginal=marginal,
        contribution=contribution,
        pct_contribution=pct_contribution,
    )
\end{lstlisting}

This is a direct transcription of the formulas above: \texttt{marginal}
is $\text{MRC}_i$, \texttt{contribution} is $\text{RC}_i$, and
\texttt{pct\_contribution} is $\%\text{RC}_i$. It is validated by three
tests in \texttt{tests/test\_contributions.py}:
\begin{itemize}
    \item \textbf{Sum invariant}: $\sum_i \text{RC}_i = \sigma_p(w)$ for
    random valid weights and a random PSD covariance matrix -- Euler's
    theorem holds numerically to \texttt{1e-8} relative tolerance.
    \item \textbf{Known two-asset case}: at the closed-form ERC weights
    of Section~\ref{sec:two-asset}, contributions are equal.
    \item \textbf{Concentrated portfolio edge case}: a 100\%-in-one-asset
    portfolio has all risk contribution on that asset and zero elsewhere.
\end{itemize}

\texttt{riskparity/optimization/erc.py::solve\_erc} (not yet implemented)
will wrap the log-barrier reformulation above to actually produce ERC
weights from a covariance matrix.

% ============================================================================
\section{Tail Risk Measures: VaR and CVaR}
\label{sec:var-cvar}
% ============================================================================

\subsection{Definitions}

For a portfolio return series $R$ (realized or simulated) and a
confidence level $\alpha$ (e.g.\ $0.975$), expressed so that both
quantities are positive loss numbers:
\[
\Var_\alpha(R) = -\,Q_{1-\alpha}(R), \qquad
\CVar_\alpha(R) = -\,\mathbb{E}\big[R \mid R \leq -\Var_\alpha(R)\big],
\]
where $Q_{1-\alpha}$ is the $(1-\alpha)$-quantile of the return
distribution. $\CVar_\alpha$ (also called Expected Shortfall) is the
average loss \emph{in the tail beyond} $\Var_\alpha$, and is a coherent
risk measure (subadditive, unlike VaR).

\subsection{Historical (non-parametric) estimator}

Directly from the empirical distribution of portfolio returns over the
lookback window: sort the realized/simulated returns, take the
$(1-\alpha)$ empirical quantile for $\Var_\alpha$, and average everything
at or below it for $\CVar_\alpha$. Simple and assumption-light, but noisy
with limited history -- with $T=500$ observations and $\alpha = 0.975$,
the CVaR estimate averages over only $\approx 12$--$13$ tail
observations.

\subsection{Cornish--Fisher (semi-parametric) estimator}

Adjusts the Gaussian quantile for sample skewness $S$ and excess kurtosis
$K$ via the Cornish--Fisher expansion of the standard normal quantile
$z_\alpha$:
\[
z_\alpha^{CF} = z_\alpha + \frac{1}{6}(z_\alpha^2 - 1)S
+ \frac{1}{24}(z_\alpha^3 - 3z_\alpha)K
- \frac{1}{36}(2z_\alpha^3 - 5 z_\alpha)S^2,
\]
then $\Var_\alpha \approx -(\mu + \sigma \, z_\alpha^{CF})$. More stable
than the historical estimator with limited data, but can misbehave for
extreme skew/kurtosis -- worth a direct comparison plot against the
historical estimator once both are implemented
(\texttt{riskparity/risk/tail\_risk.py}).

\subsection{Rockafellar--Uryasev (2000): CVaR as a convex program}

A key result used to justify \emph{minimizing} CVaR (distinct from
\emph{equalizing CVaR contributions}, which is the harder problem tackled
in Section~\ref{sec:cvar-parity}): for scenario returns $r_1, \dots, r_T$,
\[
\CVar_\alpha(w) = \min_{c \in \mathbb{R}} \left[ c + \frac{1}{(1-\alpha) T} \sum_{t=1}^T \big(-w^\top r_t - c\big)^+ \right],
\]
where $(x)^+ = \max(x, 0)$. This is convex (in fact LP-representable via
an auxiliary variable per scenario) and lets CVaR be minimized with
standard convex solvers. It does \emph{not} directly solve the
risk-parity problem -- equalizing contributions is a different, harder
objective (Section~\ref{sec:cvar-parity}).

% ============================================================================
\section{Component CVaR and CVaR-Parity}
\label{sec:cvar-parity}
% ============================================================================

\subsection{Why CVaR has no closed form in $(w, \Sigma)$}

Unlike volatility, CVaR is not a smooth quadratic function of $w$ under a
covariance matrix -- it depends on the shape of the tail of the portfolio
return distribution, which cannot in general be summarized by
second-moment information alone. It must be estimated from a scenario set
(historical returns, or simulated draws from a fitted model).

\subsection{Scenario-based marginal CVaR contribution}

Given scenario asset returns $r_{i,t}$ and portfolio returns
$R_t = w^\top r_t$, let $T_\alpha = \{t : R_t \leq -\Var_\alpha(R)\}$ be
the set of tail scenarios. The scenario-based estimator of the marginal
CVaR contribution of asset $i$ (Boudt, Peterson \& Carl, 2013, in spirit)
is
\[
\text{MRC}_i \approx -\frac{1}{|T_\alpha|} \sum_{t \in T_\alpha} r_{i,t},
\qquad
\text{RC}_i = w_i \cdot \text{MRC}_i,
\]
i.e.\ the (sign-flipped) average return of asset $i$ specifically during
the scenarios that make up the portfolio's tail. By construction
$\sum_i \text{RC}_i \approx \CVar_\alpha(w)$ -- \emph{approximately}
exact, unlike the volatility case, because it relies on the scenario tail
set being exactly the same set that defines $\CVar_\alpha(w)$ itself. The
residual numerical error should be measured and reported (planned test in
\texttt{tests/test\_contributions.py}, currently a TODO) rather than
assumed to be negligible.

\subsection{Non-convexity of CVaR-parity}

Minimizing total CVaR is convex (Section~\ref{sec:var-cvar}), but
\emph{equalizing Component CVaR contributions} is a different and
generally \textbf{non-convex} problem: the tail scenario set
$T_\alpha(w)$ itself depends on $w$ (which scenarios fall in the tail
changes as weights change), so $\text{RC}_i(w)$ is not a smooth,
well-behaved function of $w$ globally. This is the central
implementation difficulty that differentiates this project from a
superficial ``risk parity in Python'' clone, and it is treated
explicitly rather than glossed over.

\subsection{Solving algorithm: cyclical / iterative update}

Analogous in spirit to the vol-ERC algorithm, but with no guarantee of
convergence, so convergence is tracked and reported rather than assumed:

\begin{enumerate}
    \item Start from an initial guess $w^{(0)}$ (e.g.\ equal weight, or
    the vol-ERC solution).
    \item At iteration $k$: compute $\text{RC}_i(w^{(k)})$ using the
    scenario-based estimator above.
    \item Nudge weights toward equalizing contributions (e.g.\ multiply
    each $w_i$ by a factor pulling $\text{RC}_i$ toward the target budget
    $b_i \cdot \CVar_\alpha(w^{(k)})$, then renormalize to sum to 1 and
    project onto $w \geq 0$).
    \item Record the dispersion objective
    $D^{(k)} = \sum_{i,j}(\text{RC}_i(w^{(k)}) - \text{RC}_j(w^{(k)}))^2$.
    \item Stop when $D^{(k)} < \texttt{tol}$ or after
    \texttt{max\_iter} iterations (see
    \texttt{OptimizationParams} in \texttt{riskparity/config.py}).
\end{enumerate}

\texttt{riskparity/optimization/cvar\_parity.py::solve\_cvar\_parity}
(not yet implemented) must return \texttt{(weights, convergence\_history)}
so the sequence $D^{(0)}, D^{(1)}, \dots$ can be inspected and plotted --
this convergence plot is itself one of the figures in
\texttt{results/figures/} (Section~\ref{sec:results}), and if convergence
turns out to be unreliable on this universe, that is a finding to
document, not a bug to hide.

% ============================================================================
\section{Backtest Methodology}
\label{sec:backtest}
% ============================================================================

\subsection{Rolling, no-look-ahead scheme}

At each rebalance date $t$ (frequency set by
\texttt{BacktestParams.rebalance\_freq}), each strategy's weights are
computed using \emph{only} data strictly before $t$ -- the covariance
estimate or the scenario set passed to a solver must never include
information from $t$ onward. Weights are then held until the next
rebalance date, and a transaction cost of
\texttt{BacktestParams.transaction\_cost\_bps} basis points is applied to
turnover at each rebalance.

\subsection{Risk-contribution drift}

A risk-parity portfolio is only balanced \emph{at} the rebalance date --
between rebalances, prices move and weights (hence contributions) drift
away from the target. Tracking $\%\text{RC}_i$ at $t, t+1, t+2, \dots$
using the weights fixed at $t$ answers ``does the portfolio stay balanced
between rebalances,'' which matters in practice and is easy to overlook
if only the at-rebalance contributions are checked.

\subsection{Stress windows}

\begin{center}
\begin{tabular}{ll}
\toprule
Period & Window \\
\midrule
GFC 2008 & 2008-09-01 -- 2009-03-31 \\
Euro crisis 2011 & 2011-07-01 -- 2011-10-31 \\
Taper tantrum 2013 & 2013-05-01 -- 2013-09-30 \\
Vol shock 2018 & 2018-10-01 -- 2018-12-31 \\
COVID crash 2020 & 2020-02-19 -- 2020-03-23 \\
Rate shock 2022 & 2022-01-01 -- 2022-10-31 \\
\bottomrule
\end{tabular}
\end{center}
\todo[inline]{Dates are approximate placeholders from
\texttt{riskparity/backtest/stress\_periods.py}; tighten against the
actual ETF price history once data is downloaded (see
\texttt{PROJECT\_BRIEF.md}, ``Open questions'').}

% ============================================================================
\section{Comparison Metrics}
\label{sec:metrics}
% ============================================================================

Planned in \texttt{riskparity/analysis/compare.py}:
\begin{itemize}
    \item \textbf{Performance table}: annualized return, annualized
    volatility, Sharpe ratio, maximum drawdown, Calmar ratio, average
    turnover -- one row per strategy (Equal-Weight, Min-Variance,
    Vol-ERC, CVaR-Parity).
    \item \textbf{Diversification ratio}: weighted average of individual
    asset volatilities divided by portfolio volatility,
    $\text{DR}(w) = \dfrac{\sum_i w_i \sigma_i}{\sigma_p(w)} \geq 1$ --
    the classic risk-parity diversification metric.
    \item \textbf{Risk-contribution stability}: standard deviation of
    $\%\text{RC}_i$ across the backtest, per asset, per strategy -- the
    number that determines whether the CVaR-parity ``stays balanced under
    stress'' story actually holds up empirically, reported honestly
    either way.
\end{itemize}

% ============================================================================
\section{Implementation Log}
\label{sec:impl-log}
% ============================================================================

Running log, one subsection per module, updated in the same change that
implements it. Order follows the ``Suggested build order'' in
\texttt{CLAUDE.md}. Each entry should eventually contain: the formula
actually implemented (if it differs from the plan above, explain why),
a short annotated code excerpt, and what the corresponding test checks.

\subsection{\texttt{risk/volatility.py} -- sample \& shrunk covariance}

Implemented: \texttt{sample\_covariance}, \texttt{shrunk\_covariance},
\texttt{portfolio\_volatility}.

\begin{lstlisting}
TRADING_DAYS_PER_YEAR = 252

def sample_covariance(returns: pd.DataFrame, lookback_days: int) -> pd.DataFrame:
    window = returns.tail(lookback_days)
    if len(window) < 2:
        raise ValueError(f"sample_covariance: need at least 2 observations, got {len(window)}")
    daily_cov = window.cov()
    return daily_cov * TRADING_DAYS_PER_YEAR

def shrunk_covariance(returns: pd.DataFrame, lookback_days: int, shrinkage: float) -> pd.DataFrame:
    if not 0.0 <= shrinkage <= 1.0:
        raise ValueError(f"shrinkage must be in [0, 1], got {shrinkage}")
    sample_cov = sample_covariance(returns, lookback_days)
    target = pd.DataFrame(
        np.diag(np.diag(sample_cov.values)),
        index=sample_cov.index, columns=sample_cov.columns,
    )
    return (1.0 - shrinkage) * sample_cov + shrinkage * target

def portfolio_volatility(weights: np.ndarray, cov: np.ndarray) -> float:
    weights = np.asarray(weights, dtype=float)
    cov = np.asarray(cov, dtype=float)
    variance = weights @ cov @ weights
    return float(np.sqrt(max(variance, 0.0)))
\end{lstlisting}

\paragraph{Annualization.} \texttt{sample\_covariance} takes the trailing
\texttt{lookback\_days} rows of the daily log-return panel, computes the
daily sample covariance matrix, and scales it by
\texttt{TRADING\_DAYS\_PER\_YEAR = 252}. This follows from the
i.i.d.-return approximation: if daily returns were independent and
identically distributed, the variance of a $k$-day sum of returns is $k$
times the daily variance (variances of independent variables add), so a
daily covariance matrix is scaled by 252 to express the same quantity on
an annual basis. This is an approximation (returns are not exactly i.i.d.
in practice -- there is volatility clustering), but it is the standard
convention and is what makes covariance-derived quantities in this
project (e.g.\ $\sigma_p(w)$) directly comparable to commonly quoted
annualized volatility figures.

\paragraph{Shrinkage: fixed intensity toward the diagonal target, not
\texttt{sklearn}'s automatic \texttt{LedoitWolf}.} The module docstring's
TODO offered two options: a manually-chosen fixed shrinkage intensity
using \texttt{RiskParams.covariance\_shrinkage}, or calling
\texttt{sklearn.covariance.LedoitWolf}, which estimates its own
data-driven \emph{optimal} shrinkage intensity internally. These two
options are not composable: if \texttt{LedoitWolf} is used, it ignores
whatever value is set in \texttt{RiskParams.covariance\_shrinkage} and
picks its own -- which would make that config field dead, unused
configuration. Since the explicit design goal of \texttt{config.py}
(Section~\ref{sec:vol-erc-impl}, and \texttt{PROJECT\_BRIEF.md}) is that
"what happens if I change the shrinkage intensity?" should be a
one-line edit, the manual, fixed-intensity version is implemented
instead:
\[
\Sigma_{\text{shrunk}} = (1-\delta)\, \Sigma_{\text{sample}} + \delta\, T,
\qquad \delta = \texttt{RiskParams.covariance\_shrinkage},
\]
where the shrinkage target $T = \mathrm{diag}(\mathrm{diag}(\Sigma_{\text{sample}}))$
keeps every asset's variance unchanged (the diagonal) and shrinks every
off-diagonal covariance -- and therefore every implied correlation --
toward 0. The diagonal target is the specific structured target used
(rather than, e.g., a constant-correlation target) because variances are
generally much better estimated from limited history than pairwise
correlations are, so it is the correlations that benefit most from being
pulled toward a simpler, lower-variance estimate. At $\delta=0$ this
reduces exactly to the sample covariance; at $\delta=1$ it is fully
diagonal (a zero-correlation assumption across all assets).

\paragraph{Honesty about the effect size.} With only 6 assets and a
252-day (roughly one trading year) lookback window, the sample covariance
here is already reasonably well-conditioned -- shrinkage's practical
effect on this project's specific universe is expected to be modest. It
is implemented anyway because (a) it is the textbook-correct response to
estimation error that grows as the number of assets approaches the
number of observations, and demonstrating that awareness is one of the
project's explicit goals (Section~\ref{sec:qa-log} is not the place for
this, but \texttt{PROJECT\_BRIEF.md}'s "target audience" section is), and
(b) it is a one-parameter, cheaply reversible choice
(\texttt{RiskParams.covariance\_shrinkage = 0.0} disables it entirely)
that costs nothing to leave in.

\paragraph{Tests (\texttt{tests/test\_volatility.py}).}
\begin{itemize}
    \item \texttt{sample\_covariance} matches a manual
    \texttt{returns.tail(n).cov() * 252} computation, both for the
    default lookback and for a shorter one -- confirming the window
    slicing is exact, not off-by-one.
    \item Fewer than 2 observations in the window raises rather than
    returning a degenerate (all-\texttt{NaN} or divide-by-zero)
    covariance matrix.
    \item $\delta = 0$ recovers the sample covariance exactly;
    $\delta = 1$ produces an exactly diagonal matrix; at $\delta = 0.5$
    the diagonal (variances) is untouched while only the off-diagonal
    entries move -- directly testing the "shrink correlations, not
    variances" design choice above, not just the formula's algebra.
    \item An out-of-range shrinkage intensity ($\delta = 1.5$) raises.
    \item \texttt{portfolio\_volatility} matches
    $\sqrt{w^\top \Sigma w}$ computed directly, and reduces to a single
    asset's own volatility for a 100\%-concentrated portfolio -- the same
    edge case already validated for
    \texttt{volatility\_risk\_contributions} in
    \texttt{tests/test\_contributions.py}, giving two independent tests
    of the same invariant from two different modules.
\end{itemize}

\subsection{\texttt{optimization/erc.py} -- \texttt{solve\_erc}}

Implemented, via the convex log-barrier reformulation
(Section~\ref{sec:vol-erc}) solved with \texttt{cvxpy}.

\begin{lstlisting}
w = cp.Variable(n)
objective = cp.Minimize(0.5 * cp.quad_form(w, cov) - risk_budgets @ cp.log(w))
problem = cp.Problem(objective, [w >= 1e-12])
solver_tol = params.tol * 1e-2
problem.solve(solver=cp.CLARABEL, tol_gap_abs=solver_tol, tol_gap_rel=solver_tol, tol_feas=solver_tol)

raw_weights = np.asarray(w.value).flatten()
weights = raw_weights / raw_weights.sum()

achieved = volatility_risk_contributions(weights, cov).pct_contribution
dispersion = float(np.max(np.abs(achieved - risk_budgets)))
if dispersion > params.tol:
    logger.warning("... exceeding OptimizationParams.tol ...")
\end{lstlisting}

\paragraph{Why \texttt{cvxpy} over SLSQP.} The module docstring's TODO
offered a choice between SLSQP (\texttt{scipy.optimize.minimize}) on the
direct dispersion objective, or the convex log-barrier reformulation
solved with \texttt{cvxpy}. \texttt{cvxpy} was chosen: the reformulation
is provably convex with a unique minimizer (up to scale) for any strictly
positive risk-budget vector, so a convex solver either certifies it found
that global optimum or reports infeasibility -- it cannot silently settle
for a local optimum the way a general local solver could on a
hand-rolled, non-convex-looking dispersion objective. The cost is that
this choice does not reuse machinery for
\texttt{optimization/cvar\_parity.py}, which has no equivalent convex
reformulation regardless of what is chosen here (Section~\ref{sec:cvar-parity}).

\paragraph{Unequal risk budgets.} Supported directly and without extra
solver logic: the reformulation is convex for \emph{any} strictly
positive budget vector $b$ summing to 1, not only $b_i = 1/n$. Invalid
budgets (wrong length, non-positive entries, or not summing to 1) raise
\texttt{ValueError} before the solver is ever invoked.

\paragraph{An environment-specific solver-precision finding.} \texttt{cvxpy}
picks a default backend solver automatically for a problem needing the
exponential cone (required by the \texttt{cp.log} term). Left to that
default, the achieved risk-contribution dispersion after renormalizing
was only accurate to $\sim 10^{-5}$--$10^{-6}$ -- comfortably converged
(solver status \texttt{"optimal"}) but not tight enough to satisfy
\texttt{OptimizationParams.tol}$=10^{-8}$. Rather than loosening
\texttt{tol} to make the discrepancy disappear, the actual cause was
identified: explicitly requesting the \texttt{CLARABEL} backend with
tight tolerances (\texttt{tol\_gap\_abs}, \texttt{tol\_gap\_rel},
\texttt{tol\_feas}, each set to \texttt{params.tol} $\times 10^{-2}$)
brings the achieved dispersion down to $\sim 10^{-10}$, safely inside the
target tolerance. This is the same "measure, don't loosen the test until
it passes" discipline the testing philosophy calls for
(Section~\ref{sec:qa-log} has the related war story about
\texttt{cvxpy} failing to import at all on this development machine
before that was resolved).

\paragraph{Convergence verification.} After solving, the function
recomputes \%RC$_i$ at the returned weights via
\texttt{volatility\_risk\_contributions} (Section~\ref{sec:vol-erc-impl})
and compares it against the target risk budgets; a deviation larger than
\texttt{OptimizationParams.tol} is logged as a warning rather than
raised, since a numerical solver's status of \texttt{"optimal"} is not by
itself a guarantee of any particular downstream precision -- this makes
the check an observable diagnostic rather than a silent assumption.

\paragraph{Tests (\texttt{tests/test\_erc.py}).}
\begin{itemize}
    \item Weights are long-only and sum to 1.
    \item Risk contributions are equal at the returned optimum (within
    \texttt{1e-6}), fed back through
    \texttt{volatility\_risk\_contributions} rather than trusted from the
    solver's internal state.
    \item Matches the closed-form $w_i \propto 1/\sigma_i$ two-asset case
    (Section~\ref{sec:two-asset}) -- here the solver must \emph{discover}
    those weights, rather than being handed them as in
    \texttt{tests/test\_contributions.py}.
    \item A symmetric case (equal variances, equal pairwise correlation)
    gives exactly equal weights, as required by symmetry.
    \item Unequal risk budgets are honored: a 75/25 budget on the
    two-asset case produces \%RC $\approx$ (0.75, 0.25).
    \item An invalid risk-budget vector (does not sum to 1) raises
    before the solver runs.
\end{itemize}

\subsection{\texttt{risk/tail\_risk.py} -- historical \& Cornish-Fisher CVaR}
\label{sec:tail-risk-impl}

Implemented: \texttt{historical\_var}, \texttt{historical\_cvar}, and
\texttt{cornish\_fisher\_cvar}, following the formulas of
Section~\ref{sec:var-cvar}, plus a shared \texttt{\_validate\_confidence}
helper.

\begin{lstlisting}
def _validate_confidence(confidence: float) -> None:
    if not 0.0 < confidence < 1.0:
        raise ValueError(f"confidence must be in (0, 1), got {confidence}")

def historical_var(portfolio_returns: pd.Series, confidence: float) -> float:
    _validate_confidence(confidence)
    portfolio_returns = np.asarray(portfolio_returns, dtype=float)
    percentile = (1 - confidence) * 100
    var_value = np.percentile(portfolio_returns, percentile)
    return -var_value

def historical_cvar(portfolio_returns: pd.Series, confidence: float) -> float:
    _validate_confidence(confidence)
    portfolio_returns = np.asarray(portfolio_returns, dtype=float)
    percentile = (1 - confidence) * 100
    var_value = np.percentile(portfolio_returns, percentile)
    tail_losses = portfolio_returns[portfolio_returns <= var_value]
    return -np.mean(tail_losses)
\end{lstlisting}

\texttt{cornish\_fisher\_cvar} implements the expansion of
Section~\ref{sec:var-cvar} directly: sample skew $S$ and excess kurtosis
$K$ (via \texttt{scipy.stats.skew}/\texttt{kurtosis}) adjust the Gaussian
quantile $z_\alpha$ into $z_\alpha^{CF}$, and a corresponding
Cornish-Fisher expected-shortfall multiplier $e_{CF}$ (not merely the
quantile itself) is used to get \texttt{CVaR} $= -(\mu + \sigma \cdot
e_{CF})$, rather than naively averaging past observations below the
adjusted quantile -- consistent with the semi-parametric spirit of the
method (no scenario averaging at all, only the first four sample
moments).

\paragraph{Confidence validation: a real gap the tests caught.} An
earlier version of \texttt{historical\_var}/\texttt{historical\_cvar}
had no explicit check on \texttt{confidence} at all. It happened to
raise for \texttt{confidence > 1} purely as a side effect of
\texttt{np.percentile} itself rejecting a percentile argument outside
$[0, 100]$ -- but \texttt{confidence = 0.0} maps to percentile $100$,
which \texttt{np.percentile} accepts without complaint, so that
degenerate, financially meaningless confidence level silently produced
a number instead of raising. \texttt{tests/test\_tail\_risk.py::test\_invalid\_confidence\_raises}
caught exactly this asymmetry. The fix was the explicit
\texttt{\_validate\_confidence} helper above, shared by all three
functions, rather than depending on an incidental side effect of the
underlying numpy call -- the same "measure and fix the real cause, don't
paper over it" discipline already documented for \texttt{solve\_erc}'s
solver-precision finding.

\paragraph{Tests (\texttt{tests/test\_tail\_risk.py}).}
\begin{itemize}
    \item \texttt{CVaR $\geq$ VaR} for any return series with variation
    -- a structural invariant, not distribution-specific.
    \item Both increase (weakly) monotonically with the confidence
    level.
    \item A constant (zero-variation) return series reduces VaR and
    CVaR to exactly $-r$ -- the trivial edge case.
    \item On a large i.i.d.\ Gaussian sample, \texttt{historical\_var}/\texttt{historical\_cvar}
    and \texttt{cornish\_fisher\_cvar} both approximately match the
    closed-form Gaussian VaR/CVaR formula, within a documented
    \texttt{rel=0.15} tolerance -- loose because at $\alpha=0.975$ and
    $n=5000$ the CVaR tail averages only $\sim 125$ observations, so
    sampling noise on the estimate is real, not a bug to hide by
    tightening the assertion.
    \item On a synthetic negatively-skewed series (small positive drift
    plus rare large negative jumps), \texttt{cornish\_fisher\_cvar}
    reports a strictly larger CVaR than the symmetric-Gaussian
    closed-form value at the same mean/vol -- the qualitative behavior
    the whole CVaR-parity story depends on.
    \item An out-of-range \texttt{confidence} raises \texttt{ValueError}
    for both \texttt{historical\_var} and \texttt{historical\_cvar}.
\end{itemize}

\subsection{\texttt{risk/contributions.py} -- \texttt{cvar\_risk\_contributions}}
\label{sec:cvar-contrib-impl}

Implemented using the scenario-based estimator described in
Section~\ref{sec:cvar-parity}, reusing \texttt{historical\_var}
(Section~\ref{sec:tail-risk-impl}) to locate the tail rather than
re-deriving the quantile calculation a second time.

\begin{lstlisting}
def cvar_risk_contributions(weights, asset_returns, confidence) -> RiskContributions:
    weights = np.asarray(weights, dtype=float)
    asset_returns = np.asarray(asset_returns, dtype=float)

    portfolio_returns = weights @ asset_returns  # R_t, shape (n_scenarios,)
    var_alpha = historical_var(portfolio_returns, confidence)
    tail_mask = portfolio_returns <= -var_alpha
    if not np.any(tail_mask):
        raise ValueError("no scenarios fall in the tail at this confidence")

    tail_asset_returns = asset_returns[:, tail_mask]
    marginal = -np.mean(tail_asset_returns, axis=1)
    contribution = weights * marginal
    total_risk = -np.mean(portfolio_returns[tail_mask])

    return RiskContributions(
        total_risk=float(total_risk), marginal=marginal,
        contribution=contribution, pct_contribution=contribution / total_risk,
    )
\end{lstlisting}

\paragraph{The Euler sum invariant is exact here, not merely
approximate.} Section~\ref{sec:cvar-parity} notes that
$\sum_i \text{RC}_i \approx \CVar_\alpha(w)$ only \emph{approximately}
in general for the scenario estimator. In this specific implementation
it is exact: \texttt{total\_risk} and \texttt{marginal} are computed
from the \emph{same} \texttt{tail\_mask}, so
$\sum_i \text{RC}_i = \sum_i w_i \cdot \big(-\text{mean}_{t \in T}(r_{i,t})\big)
= -\text{mean}_{t \in T}\big(\sum_i w_i r_{i,t}\big)
= -\text{mean}_{t \in T}(R_t) = \text{total\_risk}$
by linearity of the mean -- confirmed numerically in
\texttt{tests/test\_contributions.py::test\_cvar\_contributions\_sum\_to\_total\_risk}
to \texttt{rel=1e-8}. The "approximate" caveat in
Section~\ref{sec:cvar-parity} would only start to bite if a different
threshold were used for the tail set than the one \texttt{historical\_var}
itself defines (e.g.\ comparing this estimator's sum against CVaR from
an independent estimator such as \texttt{cornish\_fisher\_cvar}).

\paragraph{Guarding an empty tail.} If no scenario falls at or below the
threshold (possible with a very high \texttt{confidence} or a short
scenario set), \texttt{np.mean} over an empty slice would silently
return \texttt{NaN}. This is treated as a hard \texttt{ValueError}
instead, consistent with this project's convention of surfacing
degenerate numerical situations rather than letting them propagate as
silent \texttt{NaN}s downstream (the same instinct behind
\texttt{clean\_prices} raising on an empty post-cleaning panel,
Section~\ref{sec:vol-erc-impl} is the covariance-side analogue).

\paragraph{Tests (\texttt{tests/test\_contributions.py}).}
\begin{itemize}
    \item Sum invariant: $\sum_i \text{RC}_i = \CVar_\alpha(w)$ exactly
    (to \texttt{rel=1e-8}), and $\sum_i \%\text{RC}_i = 1$.
    \item Concentrated portfolio edge case, mirroring the volatility
    version: a 100\%-in-one-asset portfolio has zero CVaR contribution
    on the untouched asset.
    \item \textbf{The central empirical claim, first evidence}: two
    assets with similar overall volatility, one symmetric and one with
    a synthetic fat left tail (same jump construction as
    Section~\ref{sec:tail-risk-impl}'s skew test), at equal 50/50
    weights. Component CVaR assigns the skewed asset a percentage risk
    contribution more than 5 percentage points higher than
    volatility-based risk contribution does for the same asset --
    i.e.\ volatility-based risk measurement is demonstrably blind to
    exactly the asymmetry Component CVaR is designed to capture. This
    is not yet the full skewed-universe divergence test at the
    portfolio-optimization level (that is
    \texttt{tests/test\_cvar\_parity.py}'s
    \texttt{test\_cvar\_parity\_differs\_from\_vol\_erc\_on\_skewed\_universe},
    still a TODO pending \texttt{solve\_cvar\_parity}), but it is the
    same phenomenon demonstrated one layer down, at the level of
    contributions rather than solved weights.
\end{itemize}

\subsection{\texttt{optimization/cvar\_parity.py} -- \texttt{solve\_cvar\_parity}}
\label{sec:cvar-parity-impl}

Implemented using the cyclical/iterative update outlined in
Section~\ref{sec:cvar-parity}: at each iteration, recompute Component
CVaR contributions (via \texttt{cvar\_risk\_contributions},
Section~\ref{sec:cvar-contrib-impl}) at the current weights, then
multiply each weight by a factor pulling its contribution toward its
budgeted target.

\begin{lstlisting}
for _ in range(params.max_iter):
    result = cvar_risk_contributions(weights, scenario_returns, confidence)
    dispersion = float(np.sum((result.pct_contribution - risk_budgets) ** 2))
    convergence_history.append(dispersion)

    if dispersion < params.tol:
        break

    adjustment_factor = 1 - params.cvar_parity_learning_rate * (
        (result.pct_contribution / risk_budgets) - 1
    )
    weights = weights * adjustment_factor
    weights = np.maximum(weights, 1e-8)
    weights = weights / weights.sum()

if convergence_history[-1] >= params.tol:
    logger.warning("solve_cvar_parity: did not converge after %d iterations ...")
\end{lstlisting}

\paragraph{The update rule.} Writing $b_i$ for the risk budget and
$\%\text{RC}_i(w^{(k)})$ for the achieved percentage Component CVaR
contribution at iteration $k$, the multiplicative factor is
\[
w_i^{(k+1)} = w_i^{(k)} \cdot \left(1 - \eta \left(\frac{\%\text{RC}_i(w^{(k)})}{b_i} - 1\right)\right),
\]
where $\eta = \texttt{OptimizationParams.cvar\_parity\_learning\_rate}$.
An asset over-contributing relative to its budget
($\%\text{RC}_i > b_i$) gets its weight shrunk; an under-contributing
one gets it grown -- the same directional intuition as Roncalli's
cyclical ERC update, just applied to Component CVaR instead of
Component volatility. Weights are floored at $10^{-8}$ (mirroring the
$w \geq 10^{-12}$ floor in \texttt{optimization/erc.py}) and
renormalized to sum to 1 after every step. Using $\%\text{RC}_i / b_i$
rather than an absolute target $b_i \cdot \CVar_\alpha(w)$ in the
denominator is a deliberate simplification: since $\%\text{RC}_i =
\text{RC}_i / \CVar_\alpha(w)$, the two are algebraically identical
($(\text{RC}_i - b_i \CVar_\alpha(w))/(b_i \CVar_\alpha(w)) =
\%\text{RC}_i/b_i - 1$), but dividing by $b_i$ -- guaranteed strictly
positive by the validation above -- avoids needing an arbitrary
$\varepsilon$ safeguard against a near-zero total-CVaR denominator.

\paragraph{Orientation convention.} \texttt{asset\_returns} is accepted
in the same $(n_{\text{scenarios}}, n_{\text{assets}})$ orientation used
everywhere else in the codebase (dates as rows, tickers as columns,
matching \texttt{clean\_prices}), and transposed once internally to the
$(n_{\text{assets}}, n_{\text{scenarios}})$ shape
\texttt{cvar\_risk\_contributions} expects -- callers should not
pre-transpose.

\paragraph{Non-convergence is reported, not hidden.} Consistent with
Section~\ref{sec:cvar-parity}'s warning that this problem is not
convex and convergence is not guaranteed, the loop does not assume
success: if the final entry of \texttt{convergence\_history} is still
$\geq \texttt{params.tol}$ after \texttt{max\_iter} iterations, a
\texttt{logger.warning} reports the final dispersion and iteration
count, symmetric to the post-hoc check already documented for
\texttt{solve\_erc} (Section~\ref{sec:vol-erc-impl}). Empirically, on
synthetic Gaussian universes with a few thousand scenarios, dispersion
drops by roughly 4 orders of magnitude (e.g.\ $\sim 1.8\times 10^{-3}
\to \sim 3\times 10^{-5}$ over 500 iterations) without fully reaching
$\texttt{tol} = 10^{-8}$ -- exactly the kind of "real but bounded"
non-convergence this project commits to reporting honestly rather than
hiding by loosening \texttt{tol} until the warning stops firing.

\paragraph{Tests (\texttt{tests/test\_cvar\_parity.py}).}
\begin{itemize}
    \item Weights are long-only and sum to 1.
    \item \texttt{convergence\_history} is finite throughout and ends
    at least 10$\times$ smaller than it started -- substantial
    improvement is required, not exact convergence or strict
    monotonicity (a strict per-step non-increase check would be too
    strong a claim for a noisy, non-convex cyclical update).
    \item Contributions are within \texttt{0.02} (2 percentage points)
    of the target risk budget at the returned weights -- deliberately
    looser than \texttt{solve\_erc}'s \texttt{1e-6}, for the reasons
    given in the test module's own docstring (scenario-estimator noise,
    not a tolerance picked to make the test pass).
    \item Unequal budgets (75/25 on a 2-asset case) are honored within
    the same tolerance.
    \item An invalid risk-budget vector raises before the loop runs.
    \item \textbf{The central empirical claim, at the solver level}:
    on a 2-asset synthetic universe with matched volatility but one
    asset carrying a strong negative skew (fat left tail, same
    construction as Section~\ref{sec:cvar-contrib-impl}'s test),
    \texttt{solve\_erc} allocates them equally (blind to skew, by
    construction), while \texttt{solve\_cvar\_parity} allocates
    meaningfully less to the skewed asset (measured $\approx 0.60/0.40$
    vs.\ vol-ERC's $\approx 0.50/0.50$ on one representative run;
    tested with a $0.05$ margin) -- Component CVaR sees tail risk that
    volatility alone cannot.
\end{itemize}

\subsection{\texttt{data/download.py} -- \texttt{download\_prices}}

Implemented. This is the only module in the codebase allowed to call
\texttt{yfinance} -- every other module downstream only ever sees a plain
\texttt{pandas.DataFrame} of prices, so the data source could later be
swapped (e.g.\ for a paid vendor) without touching any risk or
optimization code.

\begin{lstlisting}
def download_prices(universe: Universe, cache_path: str | None = "data/raw/prices_raw.parquet") -> pd.DataFrame:
    cache_file = Path(cache_path) if cache_path is not None else None
    if cache_file is not None and cache_file.exists():
        return pd.read_parquet(cache_file)

    raw = yf.download(
        list(universe.tickers),
        start=universe.start_date,
        end=universe.end_date,
        auto_adjust=True,
        progress=False,
    )
    if raw.empty:
        raise RuntimeError(f"yfinance returned no data for tickers {universe.tickers}")

    prices = raw["Close"]
    if isinstance(prices, pd.Series):
        prices = prices.to_frame(name=universe.tickers[0])
    prices = prices.reindex(columns=list(universe.tickers))

    dropped = [t for t in universe.tickers if prices[t].isna().all()]
    if dropped:
        logger.warning("Tickers with no data returned, dropped: %s", dropped)
        prices = prices.drop(columns=dropped)

    inception = prices.apply(lambda col: col.first_valid_index())
    late_starters = inception[inception > prices.index.min()]
    if not late_starters.empty:
        logger.info("Tickers with later inception dates than the universe start: %s", late_starters.to_dict())

    if cache_file is not None:
        cache_file.parent.mkdir(parents=True, exist_ok=True)
        prices.to_parquet(cache_file)

    return prices
\end{lstlisting}

Step by step, what this function actually does:

\begin{enumerate}
    \item \textbf{Cache short-circuit.} If \texttt{cache\_path} is given
    and the parquet file already exists at that path, it is loaded and
    returned immediately -- \texttt{yfinance} is never called. This is
    what makes repeated runs of \texttt{scripts/run\_pipeline.py} fast
    and usable offline, and is why \texttt{data/raw/} is gitignored
    (Section~\ref{sec:vol-erc-impl} and \texttt{.gitignore}): the cache
    is a regenerable artifact, not something to commit.

    \item \textbf{Download.} \texttt{yfinance.download} is called once
    for the full list of tickers with \texttt{auto\_adjust=True}, which
    means the returned \texttt{"Close"} column is already a
    \emph{total-return-adjusted} close (splits and dividends folded in)
    rather than the raw exchange close -- this is the series that should
    be used for return computation, since a stock split or dividend
    payment is not a real economic loss/gain and must not appear as one
    in \texttt{compute\_log\_returns} later. \texttt{progress=False}
    silences \texttt{yfinance}'s console progress bar, which is noise in
    a script/pipeline context. If the whole call returns an empty
    result (e.g.\ no network, or every ticker invalid), that is treated
    as fatal (\texttt{RuntimeError}) rather than silently propagating an
    empty DataFrame downstream.

    \item \textbf{Column extraction.} For multiple tickers,
    \texttt{yfinance.download} returns columns as a
    \texttt{(field, ticker)} \texttt{MultiIndex} (e.g.\
    \texttt{("Close", "SPY")}); indexing \texttt{raw["Close"]} collapses
    this to a plain DataFrame with one column per ticker. For the
    degenerate single-ticker case \texttt{yfinance} instead returns a
    flat frame whose \texttt{"Close"} selection is a bare
    \texttt{pd.Series} with no column name -- the
    \texttt{isinstance(prices, pd.Series)} branch re-wraps that case into
    a one-column DataFrame named after the ticker, so the function's
    return type is uniform regardless of universe size. Columns are then
    reindexed to \texttt{universe.tickers}' declared order, since
    \texttt{yfinance} does not guarantee it echoes back the input order.

    \item \textbf{Fully-missing tickers are dropped, not silently kept.}
    A ticker whose entire column is \texttt{NaN} (typo, delisted, or not
    covered by Yahoo Finance) would otherwise poison every downstream
    covariance/CVaR computation that touches it. It is detected
    (\texttt{prices[t].isna().all()}), dropped, and logged as a warning
    -- an explicit, auditable decision rather than an all-\texttt{NaN}
    column quietly surviving into \texttt{clean\_prices}.

    \item \textbf{Late-starting tickers are reported, not corrected
    here.} A ticker whose price history begins after
    \texttt{universe.start\_date} (e.g.\ a commodities or REIT ETF that
    launched later than the long-treasury ETF) is detected via
    \texttt{first\_valid\_index()} and logged at \texttt{info} level.
    Deliberately, \emph{nothing is done about it in this function} --
    deciding how to handle a partial history (drop the ticker, truncate
    the whole universe to the common start date, or something else) is a
    judgment call about the analysis, not about data acquisition, so
    it is left to \texttt{clean\_prices} (Section below). This keeps the
    module boundary from the design intent honest: \texttt{download.py}
    reports what the raw data looks like; it does not decide how to use
    it.

    \item \textbf{Cache write.} If a cache path was given (and it wasn't
    already used to short-circuit in step 1), the resulting DataFrame is
    written to parquet, creating \texttt{data/raw/} if it does not exist
    yet.

    \item \textbf{Return value.} A DataFrame indexed by trading date,
    one column per surviving ticker, values = adjusted close price --
    exactly the contract documented in the function's docstring, and the
    input type every downstream module (\texttt{clean\_prices},
    eventually \texttt{risk/volatility.py}, etc.) is written against.
\end{enumerate}

\subsection{\texttt{data/clean.py} -- \texttt{clean\_prices}, \texttt{compute\_log\_returns}}

Implemented. This module is where the late-starting/missing-data
situations that \texttt{download\_prices} only \emph{reports}
(Section above) actually get resolved into a single, fully-populated
price panel.

\begin{lstlisting}
MAX_FFILL_DAYS = 5

def clean_prices(raw_prices: pd.DataFrame) -> pd.DataFrame:
    prices = raw_prices.sort_index().ffill(limit=MAX_FFILL_DAYS)

    before = len(prices)
    prices = prices.dropna(how="any")
    after = len(prices)
    if after < before:
        logger.info(
            "clean_prices: dropped %d/%d rows still containing NaNs after "
            "forward-fill (limit=%d days) -- likely pre-inception rows for "
            "a late-starting ticker",
            before - after, before, MAX_FFILL_DAYS,
        )
    if prices.empty:
        raise ValueError("clean_prices: no rows remain after alignment/cleaning")

    return prices

def compute_log_returns(prices: pd.DataFrame) -> pd.DataFrame:
    return np.log(prices).diff().dropna(how="all")
\end{lstlisting}

\paragraph{Two-tier gap policy.} The function distinguishes two kinds of
missing data with two different fixes, rather than treating every
\texttt{NaN} the same way:
\begin{itemize}
    \item \textbf{Short gaps ($\leq$ \texttt{MAX\_FFILL\_DAYS} = 5
    trading days)} -- e.g.\ a one-off data-vendor hiccup, or a listing
    trading on a slightly different holiday calendar than the rest of
    the universe -- are \textbf{forward-filled}: the last known price is
    carried forward. This is a reasonable approximation for a short gap
    because the true (unobserved) price is very likely close to the last
    traded one.
    \item \textbf{Long gaps} -- in practice, almost always a ticker that
    simply had not started trading yet at the start of the requested
    window (e.g.\ \texttt{DBC} or \texttt{VNQ} launching after
    \texttt{Universe.start\_date = "2007-01-01"}) -- are \textbf{not}
    forward-filled indefinitely. \texttt{ffill(limit=5)} only fills up to
    5 consecutive days, so a pre-inception run of missing values longer
    than that still contains \texttt{NaN}s after the fill step. Those
    rows are then dropped by \texttt{dropna(how="any")}. The net effect
    is an \textbf{inner join on dates}: the cleaned panel starts on the
    date every surviving ticker has real data, not on
    \texttt{Universe.start\_date} -- fabricating five years of a
    constant "first price" for a late-starting ETF would silently
    understate its true volatility and bias every risk measure computed
    on it, which is worse than simply shortening the usable history.
    This is logged at \texttt{info} level so the size of the truncation
    is visible, not hidden inside a silently shorter DataFrame.
    \item If cleaning leaves \emph{no} rows at all (e.g.\ two tickers
    with non-overlapping histories), that is treated as a fatal
    \texttt{ValueError} rather than returning an empty, technically
    valid but useless DataFrame to the rest of the pipeline.
\end{itemize}

\paragraph{Log returns.} $r_t = \ln(P_t / P_{t-1})$ is used everywhere in
this project instead of simple returns
$\tilde r_t = P_t/P_{t-1} - 1$, for two reasons relevant to the risk
engines built on top of it:
\begin{itemize}
    \item \textbf{Time-additivity.} The $k$-day log return is simply the
    sum of the daily log returns,
    $\ln(P_t/P_{t-k}) = \sum_{j=0}^{k-1} r_{t-j}$, which is what makes
    rolling-window aggregation (covariance over a lookback window, CVaR
    over a scenario window) straightforward. Simple returns do not add
    this way across time.
    \item \textbf{Symmetry.} Log returns treat a doubling and a halving
    of price symmetrically ($\ln 2 = -\ln(0.5)$), which is a better
    behaved input to variance/tail estimation than simple returns, whose
    distribution is skewed purely as an artifact of the compounding
    definition even if the underlying price process were symmetric.
\end{itemize}
\texttt{compute\_log\_returns} takes $\ln$ of the whole price panel
(vectorized, one call for every ticker) and then \texttt{diff()}s along
the date axis; the resulting first row is always \texttt{NaN} (there is
no previous price to compare the first observation to) and is dropped
with \texttt{dropna(how="all")}.

\paragraph{Tests (\texttt{tests/test\_data\_clean.py}).}
\begin{itemize}
    \item A single-day \texttt{NaN} is forward-filled and the row
    survives.
    \item A 7-day leading \texttt{NaN} run (longer than
    \texttt{MAX\_FFILL\_DAYS}) is \emph{not} filled -- those rows are
    dropped and the cleaned panel starts on the first date where every
    ticker has real data.
    \item An all-\texttt{NaN} ticker causes \texttt{clean\_prices} to
    raise, rather than silently returning an empty frame.
    \item \texttt{compute\_log\_returns} matches a direct
    \texttt{np.log(prices).diff()} computation, and is exactly zero every
    day for a constant (no-movement) price series -- a trivial but
    useful sanity check that the sign and scale of the transform are
    right.
\end{itemize}

\subsection{\texttt{backtest/engine.py} -- \texttt{run\_backtest}}
\label{sec:backtest-impl}

Implemented as the rolling, no-look-ahead scheme described in
Section~\ref{sec:backtest}. \texttt{run\_backtest} takes a black-box
\texttt{strategy\_fn: window\_returns -> weights} (a closure the caller
builds around \texttt{solve\_erc}, \texttt{solve\_cvar\_parity}, or
anything else), so the engine itself never imports an optimizer and
stays agnostic to which strategy it is backtesting.

\begin{lstlisting}
for i, rebal_date in enumerate(rebalance_dates):
    loc = dates.get_loc(rebal_date)
    window = returns.iloc[:loc]  # strictly before rebal_date

    if len(window) < RISK_PARAMS.vol_lookback_days:
        continue  # not enough history yet

    target_weights = np.asarray(strategy_fn(window), dtype=float)
    turnover = float(np.abs(target_weights - current_weights).sum())
    ...
    held_weights = target_weights.copy()
    for t_loc in range(loc, next_loc):
        day_returns = returns.iloc[t_loc].to_numpy()
        gross_return = float(held_weights @ day_returns)
        cost = turnover * params.transaction_cost_bps / 1e4 if t_loc == loc else 0.0
        daily_returns.iloc[t_loc] = gross_return - cost
        held_weights = held_weights * (1.0 + day_returns)
        held_weights = held_weights / held_weights.sum()
    current_weights = held_weights
\end{lstlisting}

\paragraph{No-look-ahead.} Enforced structurally, not just by
convention: \texttt{window = returns.iloc[:loc]} slices by integer
position, so it excludes the return realized on \texttt{rebal\_date}
itself (and everything after) by construction -- \texttt{strategy\_fn}
cannot see the return it is about to be judged on.

\paragraph{Rebalance dates.} Computed by
\texttt{index.to\_series().resample(freq).last().dropna()}: the last
trading date actually present in \texttt{returns} within each
\texttt{freq} period (e.g.\ the last trading day of each month for
\texttt{freq="ME"}), rather than a fixed calendar grid -- a period with
no trading days in \texttt{returns} (a holiday-heavy stretch, or a gap
in the data) is dropped instead of producing a spurious rebalance date
at a \texttt{NaT}.

\paragraph{Warm-up.} A rebalance is skipped until at least
\texttt{RiskParams.vol\_lookback\_days} (252) days of history are
available, since the ex-ante risk-contribution diagnostic below needs a
non-degenerate covariance estimate. If \emph{no} rebalance date ever
clears that bar (\texttt{returns} shorter than one lookback window),
\texttt{run\_backtest} raises \texttt{ValueError} rather than returning
a silently empty or degenerate result.

\paragraph{Weight drift between rebalances.} Weights are not re-targeted
daily: \texttt{held\_weights} drifts with realized returns
(buy-and-hold) between rebalance dates,
\[
w_i(t+1) = \frac{w_i(t)\,(1 + r_i(t))}{\sum_j w_j(t)\,(1 + r_j(t))},
\]
renormalized each day so weights keep summing to 1. This is what makes
turnover -- and, in principle, risk-contribution drift
(Section~\ref{sec:backtest}) -- a non-trivial quantity instead of zero
by construction.

\paragraph{Transaction costs.} Charged once per rebalance, proportional
to L1 turnover $\sum_i |w_i^{\text{target}} - w_i^{\text{held, just
before}}|$, deducted from that rebalance day's return only (every other
day in the holding period is unaffected). The very first rebalance is
charged against an implicit all-cash starting position
(\texttt{current\_weights = 0}), so the cost of funding the portfolio
initially is included rather than treated as free.

\paragraph{Risk-contribution lens: volatility, for every strategy.}
\texttt{risk\_contributions\_history} is computed with
\texttt{volatility\_risk\_contributions} (Section~\ref{sec:vol-erc-impl})
at every rebalance date, \emph{even for a CVaR-parity strategy}. This is
deliberate, not an oversight: using one common lens for every strategy
is what makes it possible to ask this project's central question inside
the backtest, over time, rather than only at a single static date --
does a portfolio built to equalize Component CVaR end up with unequal
\emph{volatility} contributions in practice, and does that gap widen
under stress? A CVaR-side version of the same diagnostic is a natural
extension, not implemented here: the engine has no reliable way to know
what scenario window a black-box \texttt{strategy\_fn} would consider
appropriate for its own CVaR estimate.

\paragraph{Known gap: daily risk-contribution drift is not exposed.}
Section~\ref{sec:backtest}'s ``Risk-contribution drift'' subsection
describes tracking $\%\text{RC}_i$ at every day between rebalances,
using the weights fixed at the last rebalance -- \texttt{run\_backtest}
computes the equivalent \texttt{held\_weights} internally, day by day
(needed to get \texttt{daily\_returns} right), but currently only
persists weights and risk contributions \emph{at} rebalance dates
(\texttt{weights\_history}, \texttt{risk\_contributions\_history}), not
the intermediate daily values. Reported honestly here rather than
silently narrowing the methodology section to match the code: exposing
this would mean adding a daily-frequency field to
\texttt{BacktestResult}, deferred until the comparison/plotting stage
(Section~\ref{sec:metrics}) actually needs it.

\paragraph{Tests (\texttt{tests/test\_backtest.py}).}
\begin{itemize}
    \item \texttt{\_rebalance\_dates} picks the last trading day of each
    period and drops periods with no trading days at all (not a
    \texttt{NaT}-derived spurious date).
    \item Too little history raises \texttt{ValueError}; rebalance dates
    before \texttt{RiskParams.vol\_lookback\_days} of history is
    available are skipped, checked against \texttt{\_rebalance\_dates}'s
    own output directly.
    \item \texttt{weights\_history} rows are long-only and sum to 1;
    \texttt{risk\_contributions\_history} rows sum to 1.
    \item \textbf{No-look-ahead}, checked directly: a spy
    \texttt{strategy\_fn} records every window it receives, and each one
    is asserted to end strictly before its rebalance date.
    \item Turnover on the first rebalance equals 1.0 exactly (funding
    from an all-cash start); the realized return on a rebalance day
    matches the exact formula $w \cdot r_t - \text{turnover} \times
    \text{bps}/10^4$ (no drift has occurred yet that day); comparing a
    0~bps run against a 50~bps run isolates transaction costs to
    rebalance days only, with identical returns everywhere else.
\end{itemize}

\subsection{\texttt{analysis/compare.py} -- \texttt{performance\_table}, \texttt{risk\_contribution\_stability}}
\label{sec:compare-impl}

Both functions operate purely on the \texttt{BacktestResult}s already
produced by \texttt{run\_backtest} (Section~\ref{sec:backtest-impl}) --
no re-fetching of prices or re-solving of any optimizer here, keeping
this module a thin aggregation layer over backtest output.

\paragraph{\texttt{performance\_table}.} One row per strategy:
annualized return, annualized volatility, Sharpe ratio, max drawdown,
Calmar ratio, average turnover.
\[
\text{ann. return} = (1+R_{\text{cum}})^{252/n} - 1, \qquad
\text{ann. vol} = \hat\sigma(r) \sqrt{252}, \qquad
\text{Sharpe} = \frac{\text{ann. return}}{\text{ann. vol}},
\]
\[
\text{max DD} = -\min_t\left(\frac{W_t}{\max_{s \leq t} W_s} - 1\right), \qquad
\text{Calmar} = \frac{\text{ann. return}}{\text{max DD}},
\]
where $W_t = \prod_{s \leq t}(1+r_s)$ is the wealth curve built from
\texttt{BacktestResult.returns}. Sharpe assumes a 0\% risk-free rate:
this project only sources free ETF price data (no risk-free/cash-rate
series), so introducing a nonzero rate would mean adding a data
dependency purely for this one ratio -- a deliberate simplification,
not an oversight. Max drawdown and Calmar's denominator are reported as
\emph{positive} loss magnitudes, matching this project's
risk-as-positive-number convention (Section~\ref{sec:tail-risk-impl}'s
CVaR sign flip), not the negative-number convention some libraries use.

A strategy with genuinely zero variance or drawdown (e.g.\ constant
daily returns, only possible with synthetic test data) does not land on
exactly \texttt{0.0} for the denominator -- \texttt{float64} summation
noise leaves a residual on the order of $10^{-19}$, which is still
technically positive but not a meaningful denominator. Both ratios are
therefore guarded with an explicit $\varepsilon = 10^{-12}$ threshold
(matching the numerical floor already used in
\texttt{optimization/erc.py}, Section~\ref{sec:vol-erc-impl}) rather
than an exact \texttt{> 0} comparison, returning \texttt{NaN} below it
instead of an enormous, meaningless ratio -- caught by this module's own
test suite before being fixed.

\paragraph{\texttt{risk\_contribution\_stability}.} Per strategy, per
asset, the standard deviation of $\%\text{RC}_i$ across every rebalance
date in the backtest (\texttt{BacktestResult.risk\_contributions\_history},
Section~\ref{sec:backtest-impl}). This is the number that determines
whether a strategy's ``stays balanced over time'' story
(Section~\ref{sec:backtest}, ``Risk-contribution drift'') actually
holds up: a low std means $\%\text{RC}_i$ barely moved from its target
across the whole backtest; a high one means the portfolio drifted
meaningfully out of balance between rebalances, however good it looked
at any single rebalance date.

\paragraph{Deferred: \texttt{diversification\_ratio}.} Also listed in
the original module TODO and in Section~\ref{sec:metrics}, but
deliberately not implemented in this pass: $\text{DR}(w) = \sum_i w_i
\sigma_i / \sigma_p(w)$ needs each asset's \emph{individual} volatility
and the portfolio volatility at each rebalance date, and
\texttt{BacktestResult} currently only stores percentage risk
\emph{contributions} ($\text{RC}_i/\sigma_p$), not the underlying
$\sigma_i$ or the covariance matrix used to compute them. Implementing
it properly would mean extending \texttt{BacktestResult}'s schema in
\texttt{backtest/engine.py} -- a bigger change than filling in this
module -- so it is flagged here rather than worked around with a silent
re-derivation that could drift out of sync with the engine's actual
covariance estimate.

\paragraph{Tests (\texttt{tests/test\_compare.py}).}
\begin{itemize}
    \item \texttt{performance\_table}: with exactly 252 daily
    observations the annualization exponent is 1, so annualized return
    must equal raw cumulative return exactly; a hand-constructed
    flat/-10\%/flat return path gives an exactly-checkable max drawdown
    of 0.10; average turnover is checked against a plain \texttt{mean()}
    on a hand-picked turnover series; zero-volatility returns give
    \texttt{NaN} Sharpe and Calmar (the bug described above, now
    covered by a regression test); one row per strategy with the
    expected column set.
    \item \texttt{risk\_contribution\_stability}: matches
    \texttt{pandas}' own \texttt{std(ddof=1)} on a hand-built
    contributions history exactly; a strategy with $\%\text{RC}_i$
    swinging between 0.30 and 0.70 shows a strictly higher std than one
    swinging between 0.49 and 0.51, per asset; index/columns are
    strategy names and tickers respectively.
\end{itemize}

\subsection{\texttt{analysis/plots.py}}
\label{sec:plots-impl}

Three matplotlib helpers, one per figure from Section~\ref{sec:metrics}'s
list that maps directly onto a single \texttt{BacktestResult}-derived
input; all three share a \texttt{\_finalize(fig, save\_path)} helper
that always returns the \texttt{Figure} (so a notebook cell ending in a
call to one of these renders it inline), additionally writes it to
\texttt{save\_path} if one is given, and \texttt{plt.close(fig)}s it
afterwards purely to keep matplotlib's internal figure registry from
growing unbounded across a pipeline run -- the returned object stays
fully usable.

\paragraph{\texttt{plot\_risk\_contributions}.} Grouped bar chart of
$\%\text{RC}_i$ per asset, one group of bars per strategy, at a single
date, with a dashed reference line at $1/N_{\text{assets}}$ (the
default equal risk budget used throughout this project).

\paragraph{\texttt{plot\_cumulative\_returns}.} One line per strategy,
built from \texttt{BacktestResult.returns} as $W_t =
\prod_{s\leq t}(1+r_s)$ -- the same wealth curve
\texttt{performance\_table} (Section~\ref{sec:compare-impl}) uses
internally for drawdown.

\paragraph{\texttt{plot\_contribution\_stability}.} One subplot per
strategy, stacked vertically with a shared x-axis, each showing every
asset's $\%\text{RC}_i$ across all rebalance dates
(\texttt{risk\_contributions\_history}), with the same $1/N$ reference
line -- the visual counterpart of
\texttt{risk\_contribution\_stability}'s std-dev numbers
(Section~\ref{sec:compare-impl}).

\paragraph{Deferred: a dedicated stress-window figure.} Section~\ref{sec:metrics}
also lists a fourth figure, zooming into 1--2 stress windows from
\texttt{riskparity.backtest.stress\_periods.STRESS\_PERIODS} to show
where vol-ERC and CVaR-parity diverge most. No separate function is
implemented for it: it is the exact same chart as
\texttt{plot\_cumulative\_returns}, called with each series in
\texttt{returns} pre-sliced to a stress window's date range (e.g.\
\texttt{\{name: r.loc[start:end] for name, r in returns.items()\}}) --
a second function differing only in what slice of data it receives
would duplicate \texttt{plot\_cumulative\_returns} for no reason.

\paragraph{Consistent color per strategy across figures.} Achieved via
matplotlib's default color cycle (colors assigned to series in
iteration order), not a hard-coded strategy-name-to-color mapping --
this module has no reason to know what strategy names a caller
compares. This holds automatically as long as the caller passes
strategies in the same key order to every plotting function (a plain
Python \texttt{dict} preserves insertion order), the natural way
\texttt{scripts/run\_pipeline.py} would build its \texttt{results} dict
once and reuse it across all three calls.

\paragraph{Manual sanity check.} Run against a synthetic 4-asset,
$\sim$3-year universe with \texttt{solve\_erc} and
\texttt{solve\_cvar\_parity} both wired through \texttt{run\_backtest}:
since the synthetic assets were drawn i.i.d.\ with identical
distribution parameters, vol-ERC's $\%\text{RC}_i$ stayed flat at
$0.25$ to within $\sim 10^{-6}$ across every rebalance in
\texttt{plot\_contribution\_stability}'s top subplot, while
CVaR-parity's swung between roughly $0.16$ and $0.32$ in the bottom
subplot -- visually consistent with the
\texttt{solve\_cvar\_parity: did not converge} warnings
(Section~\ref{sec:cvar-parity-impl}) logged during the same run, and
with \texttt{risk\_contribution\_stability} reporting
CVaR-parity's std roughly $10^5\times$ vol-ERC's on that run. Exactly
the kind of honest, real instability this project commits to surfacing
rather than hiding.

\paragraph{Tests (\texttt{tests/test\_plots.py}).}
\begin{itemize}
    \item Structural checks only (axis/line/bar counts, titles, whether
    a file is actually written to a not-yet-existing directory) --
    deliberately not pixel/image comparison, which would be brittle and
    add no real coverage for this module.
    \item \texttt{plot\_risk\_contributions}: exactly $N_{\text{strategies}}
    \times N_{\text{assets}}$ bars plus the one $1/N$ reference line, at
    the expected height.
    \item \texttt{plot\_cumulative\_returns}: one line per strategy; the
    final data point of a constant-return series matches
    $(1+r)^n$ exactly.
    \item \texttt{plot\_contribution\_stability}: one subplot per
    strategy, titled with the strategy name, each containing
    $N_{\text{assets}} + 1$ lines (assets plus the reference line).
    \item All three: \texttt{save\_path} produces a file on disk,
    including creating a not-yet-existing parent directory.
\end{itemize}

% ============================================================================
\section{Empirical Results}
\label{sec:results}
% ============================================================================

\todo[inline]{Populate once \texttt{scripts/run\_pipeline.py} produces
output in \texttt{results/}. Figures below are wired to appear
automatically (via \texttt{IfFileExists}) once the corresponding file is
generated -- no need to edit the includegraphics calls later, just run
the pipeline and recompile.}

\paragraph{Risk contributions, single date (Vol-ERC vs.\ CVaR-Parity).}
\IfFileExists{../results/figures/risk_contributions.png}{
    \begin{center}\includegraphics[width=0.8\textwidth]{../results/figures/risk_contributions.png}\end{center}
}{\todo[inline]{Figure not yet generated: \texttt{results/figures/risk\_contributions.png}}}

\paragraph{Cumulative returns, all strategies.}
\IfFileExists{../results/figures/cumulative_returns.png}{
    \begin{center}\includegraphics[width=0.8\textwidth]{../results/figures/cumulative_returns.png}\end{center}
}{\todo[inline]{Figure not yet generated: \texttt{results/figures/cumulative\_returns.png}}}

\paragraph{Rolling risk-contribution stability.}
\IfFileExists{../results/figures/contribution_stability.png}{
    \begin{center}\includegraphics[width=0.8\textwidth]{../results/figures/contribution_stability.png}\end{center}
}{\todo[inline]{Figure not yet generated: \texttt{results/figures/contribution\_stability.png}}}

\paragraph{CVaR-parity convergence diagnostic.}
\IfFileExists{../results/figures/cvar_parity_convergence.png}{
    \begin{center}\includegraphics[width=0.8\textwidth]{../results/figures/cvar_parity_convergence.png}\end{center}
}{\todo[inline]{Figure not yet generated: \texttt{results/figures/cvar\_parity\_convergence.png}}}

\paragraph{Performance table.}
\todo[inline]{Insert \texttt{results/tables/performance\_table} once
\texttt{analysis/compare.py::performance\_table} is implemented.}

% ============================================================================
\section{Session Q\&A Log}
\label{sec:qa-log}
% ============================================================================

An append-only log of questions asked during work sessions on this
project, with the answers given, kept for the same reason the rest of
this notebook exists: so reasoning and context survive past the chat
window they happened in. Entries are added chronologically, translated
to English if asked in another language, and are not removed or
rewritten once logged (a later entry may supersede an earlier one, but
the earlier one stays as a record of how understanding evolved).

\begin{description}[style=nextline,leftmargin=0pt]

\item[Q: Does Claude Code retain memory of previous conversations on this project?]
No. Each session starts without memory of prior chat sessions; nothing
persists automatically. The only things that carry over between sessions
are whatever gets written to durable files in the repository -- which is
exactly why \texttt{CLAUDE.md}, \texttt{PROJECT\_BRIEF.md}, and this
notebook exist: they are the persistence layer for context that would
otherwise be lost.

\item[Q: What is the purpose of a \texttt{CLAUDE.md} file?]
It is a project-level context file that Claude Code reads automatically
at the start of a session in a repository. It holds repo-specific
working instructions -- common commands (build/lint/test), architecture,
conventions, and known pitfalls -- so they do not need to be repeated in
every conversation. It can live at the repo root (shared with a team via
git) and/or globally in the user's home directory for personal
cross-project preferences.

\item[Q: What is the difference between \texttt{PROJECT\_BRIEF.md} and \texttt{CLAUDE.md}?]
They serve different audiences and purposes rather than duplicating one
another. \texttt{PROJECT\_BRIEF.md} is a design log written for a human
(the project owner, later, or a recruiter/interviewer): it records
\emph{why} the project is scoped and structured the way it is, and which
alternatives were considered and rejected. \texttt{CLAUDE.md} is
operational guidance written for the AI assistant: \emph{how} to work in
the repository -- commands to run, the architecture map, invariants not
to break. This notebook, \texttt{docs/notes.tex}, complements both as the
place where the actual finance/math/code \emph{content} lives, updated as
the project advances.

\item[Q: Can we already test whether \texttt{python scripts/run\_pipeline.py} runs end to end?]
At the time of asking: no. Nearly every function in \texttt{riskparity/}
was still a stub (\texttt{raise NotImplementedError}); only
\texttt{volatility\_risk\_contributions} was implemented. Running the
pipeline first failed on a packaging problem --
\texttt{ModuleNotFoundError: No module named 'riskparity'} -- because the
package had no \texttt{pyproject.toml} and was never installed into the
project's \texttt{.venv}. Adding a minimal \texttt{pyproject.toml} and
running \texttt{pip install -r requirements.txt -e .} fixed that; after
the fix, the pipeline correctly progressed to and failed at the first
real stub, \texttt{download\_prices}, which is the expected state of an
unfinished implementation, not a plumbing bug.

\item[Q: Where do I add Python libraries, and how?]
In \texttt{requirements.txt} at the repository root -- that is the single
place dependencies are declared. To install what is already listed:
\texttt{pip install -r requirements.txt} inside the project's
\texttt{.venv}. To add a new dependency: append a line to
\texttt{requirements.txt} (ideally with a minimum version constraint,
matching the style of the existing entries) and re-run the same install
command; pip only needs to install the new line since the rest are
already satisfied.

\item[Q: I don't see the result of the assertions in your tests -- how do I know what the code actually produced?]
This is expected \texttt{pytest} behavior, not a gap in the tests:
\texttt{pytest} only reports PASSED/FAILED by design and does not print
intermediate values for a passing assertion -- values are only shown
automatically when an assertion \emph{fails}, via \texttt{pytest}'s
assertion-rewriting diff. To see the actual computed values on a passing
run, either add explicit \texttt{print()} calls and run with
\texttt{pytest -s}, or run a short standalone script that calls the
function directly and prints its output (done for
\texttt{clean\_prices}/\texttt{compute\_log\_returns} to confirm the
forward-fill and truncation behavior visually, e.g.\ the 7-day
pre-inception gap for a late-starting ticker being dropped rather than
filled).

\item[Q: \texttt{cvxpy} fails to import with a DLL load error blocked by "a stratégie de contrôle d'application" (an application control policy) -- what is going on, and how was it resolved?]
The failure was traced to a native compiled dependency of \texttt{cvxpy}
(\texttt{qdldl}, imported eagerly by \texttt{cvxpy}'s own
\texttt{\_\_init\_\_} regardless of which solver backend is actually
used) being blocked by \textbf{Windows Smart App Control}, a system-wide
policy on this development machine that only allows binaries with a
recognized reputation/signature to load
(\texttt{VerifiedAndReputablePolicyState = 1} in the registry). Unlike an
ordinary "downloaded file" block, this is not fixed by
\texttt{Unblock-File} (which only clears the NTFS zone-identifier
stream, not an application-control decision) -- Microsoft's own
documentation states Smart App Control normally cannot be turned off
once enabled without reinstalling Windows. In this case the user
disabled it directly; after that, \texttt{import cvxpy} succeeded and
\texttt{solve\_erc}'s existing \texttt{cvxpy}-based implementation
(Section~\ref{sec:impl-log}) worked without needing to be rewritten
around a different solver library.

\item[Q: Explain in detail the risk-contribution dispersion check at the end of \texttt{solve\_erc} (\texttt{riskparity/optimization/erc.py}, the block recomputing \texttt{achieved} and comparing it to \texttt{risk\_budgets}).]
This block is a post-hoc numerical verification, not a step needed to
produce the answer. After the convex solve, \texttt{weights} are
renormalized (line~92); the block then recomputes the actually-achieved
percentage risk contributions by calling
\texttt{volatility\_risk\_contributions(weights, cov).pct\_contribution}
-- the same reference function documented in
Section~\ref{sec:vol-erc-impl} -- rather than trusting the solver's
internal state. \texttt{dispersion} is the largest absolute deviation
(infinity norm) between the achieved contributions and the target
\texttt{risk\_budgets}. If it exceeds \texttt{OptimizationParams.tol}, a
\texttt{logger.warning} records the magnitude of the deviation, the
tolerance it exceeded, and the solver's own status string
(\texttt{"optimal"} vs.\ \texttt{"optimal\_inaccurate"}) for diagnosis.
It is a warning, not a raised exception, deliberately: the convex
reformulation's optimum is exact in theory, but a numerical solver only
reaches it to floating-point/solver tolerance, so treating a small
overshoot as fatal would break the pipeline for cases that are still
perfectly usable -- consistent with this project's rule (\texttt{CLAUDE.md})
of reporting numerical imprecision honestly via logs rather than hiding
it \emph{or} treating it as an unconditional failure.

\item[Q: Why does \texttt{solve\_erc} accept solver status \texttt{"optimal\_inaccurate"} as valid, alongside \texttt{"optimal"} (line~88)?]
\texttt{"optimal\_inaccurate"} is a legitimate CLARABEL/\texttt{cvxpy}
convergence outcome, not a failure: it means the solver reached a point
consistent with the optimum but did not fully satisfy the requested
convergence tolerances (\texttt{tol\_gap\_abs}, \texttt{tol\_gap\_rel},
\texttt{tol\_feas}, set on line~86 to
\texttt{params.tol}$\times 10^{-2}$ -- deliberately tighter than the
tolerance ultimately required). Genuine failures
(\texttt{"infeasible"}, \texttt{"unbounded"}, etc.) are distinct
statuses and are excluded by the same check. Accepting
\texttt{"optimal\_inaccurate"} here does not mean precision is
unchecked: the dispersion block described in the previous entry
recomputes the actual achieved contributions regardless of which of
the two accepted statuses was returned, so the solver's status string
is only used to reject hard failures, while the real precision
judgment is made numerically on the recomputed result, not trusted
from the label alone.

\item[Q: Explain \texttt{riskparity/risk/contributions.py} in detail.]
Covered in full in Section~\ref{sec:vol-erc-impl} (the
\texttt{RiskContributions} dataclass and the implemented
\texttt{volatility\_risk\_contributions}) and
Section~\ref{sec:cvar-parity} (the stubbed
\texttt{cvar\_risk\_contributions}, not yet implemented). Separately, a
stray file \texttt{riskparity/risk/contrib.py} was found alongside it
during this session -- not referenced anywhere in \texttt{CLAUDE.md} or
imported by any other module, and containing broken code: a
non-existent type (\texttt{np.darray}, a typo for \texttt{np.ndarray})
and a function (\texttt{volatility\_risk\_contrib}) defined nested
inside the \texttt{RiskContrib} dataclass body without \texttt{self} or
a \texttt{@staticmethod} decorator, which would not work if called. It
appears to be an abandoned duplicate draft of this module, similar in
kind to a stray \texttt{ercc.py} found next to
\texttt{optimization/erc.py} earlier in the session -- flagged to the
project owner rather than deleted unilaterally.

\item[Q: What is \texttt{@dataclass}, and is it needed for \texttt{RiskContributions}?]
\texttt{@dataclass} (from the standard-library \texttt{dataclasses}
module) generates \texttt{\_\_init\_\_}, \texttt{\_\_repr\_\_}, and
\texttt{\_\_eq\_\_} automatically from a class's type-annotated fields,
instead of those being written by hand. It is a good fit for
\texttt{RiskContributions} (Section~\ref{sec:vol-erc-impl}) because
that class is a pure result container -- four arrays/scalars that
always travel together with no behavior of their own -- which is
exactly the use case the decorator targets; it is what allows the
keyword-argument construction
(\texttt{RiskContributions(total\_risk=..., marginal=..., ...)}) used
in \texttt{contributions.py}. The stray \texttt{contrib.py} file
(previous entry) shows the decorator's actual limit: it only
auto-generates behavior for declared fields, so a method carelessly
nested inside the dataclass body without \texttt{self} or
\texttt{@staticmethod} is still broken -- \texttt{@dataclass} does not
protect against that kind of structural mistake.

\item[Q: What is the current state of the project, and what should we work on next?]
At the start of this session, \texttt{optimization/cvar\_parity.py}'s
\texttt{solve\_cvar\_parity} was already implemented and staged (see
Section~\ref{sec:cvar-parity-impl}). \texttt{backtest/engine.py} was
mid-draft: \texttt{run\_backtest} was still a stub
(\texttt{print(); return 0}), and the file additionally contained
module-level exploratory code -- an inline \texttt{yfinance} download of
\texttt{AGG/SPY/EEM/BTC-USD} that executed on import -- which violates
the project rule that only \texttt{data/download.py} may talk to
\texttt{yfinance} (\texttt{CLAUDE.md}, ``Repository structure''). Agreed
next step: implement \texttt{run\_backtest} properly and remove the
exploratory code, which Section~\ref{sec:backtest-impl} now documents.

\item[Q: Implement \texttt{run\_backtest}, and explain every design choice made along the way.]
Covered in full in Section~\ref{sec:backtest-impl}: no-look-ahead via
integer-position slicing before the rebalance date, rebalance dates via
\texttt{resample(freq).last().dropna()}, a warm-up threshold of
\texttt{RiskParams.vol\_lookback\_days} before the first rebalance is
attempted, buy-and-hold weight drift between rebalances, turnover-based
transaction costs charged once per rebalance against an implicit
all-cash starting position, and the deliberate choice to report
ex-ante risk contributions with \texttt{volatility\_risk\_contributions}
for every strategy -- including CVaR-parity ones -- so the CVaR-vs-vol
divergence question can be tracked over the whole backtest, not just at
one static date.

\item[Q: Explain \texttt{run\_backtest} line by line.]
Walked through in the session: \texttt{\_rebalance\_dates} converts the
returns index to a \texttt{Series} of dates on itself purely so
\texttt{.resample().last()} can be used, keeping the last trading date
per period and dropping empty periods via \texttt{.dropna()}. The main
loop, for each rebalance date, slices \texttt{window = returns.iloc[:loc]}
(strictly before that date), skips the date if
\texttt{RiskParams.vol\_lookback\_days} of history is not yet available,
calls \texttt{strategy\_fn(window)} for the target weights, computes
turnover against the currently-held weights, records ex-ante risk
contributions on that same window, then runs an inner day-by-day loop
from the rebalance date to the next one that (a) computes that day's
gross portfolio return as \texttt{held\_weights @ day\_returns}, (b)
subtracts the transaction cost only on the rebalance day itself, and
(c) drifts \texttt{held\_weights} forward with that day's realized
returns for the next iteration. After the loop, results accumulated in
plain \texttt{dict}s are assembled into the \texttt{BacktestResult}
dataclass (Section~\ref{sec:backtest-impl}).

\item[Q: What exactly does \texttt{run\_backtest} return?]
A \texttt{BacktestResult} with four fields: \texttt{returns} (daily net
portfolio returns, from the first successful rebalance onward);
\texttt{weights\_history} (target weights, one row per rebalance date,
summing to 1); \texttt{turnover} (one scalar per rebalance date); and
\texttt{risk\_contributions\_history} (percentage volatility risk
contributions, one row per rebalance date, summing to 1). Note what is
\emph{not} returned: the day-by-day drifted weights between rebalance
dates are computed internally (\texttt{held\_weights}) but discarded
after each day, not exposed -- see the ``known gap'' paragraph in
Section~\ref{sec:backtest-impl}.

\item[Q: What does ``weight drift'' mean?]
Between two rebalance dates, nothing is bought or sold, but the assets
that outperform mechanically grow to a larger share of the portfolio
and the underperformers shrink -- e.g.\ starting at 50/50, one asset up
2\% and the other down 1\% on the same day moves the portfolio to
roughly 50.75/49.25 with zero trading. \texttt{held\_weights} in
\texttt{run\_backtest}'s inner loop (Section~\ref{sec:backtest-impl})
tracks exactly this, day by day, and it is what the next rebalance's
turnover is measured against (the actual currently-held mix, not last
month's target).

\item[Q: What is the formula used to renormalize drifted weights?]
\[
w_i(t+1) = \frac{w_i(t)\,(1 + r_i(t))}{\sum_j w_j(t)\,(1 + r_j(t))}.
\]
The numerator revalues each asset's holding by its own return; the
denominator (summed over all assets) restores the total to 1 so the
result is again expressed as portfolio-fraction weights rather than raw
revalued amounts. Implemented as \texttt{held\_weights *= (1 +
day\_returns)} followed by \texttt{held\_weights /=
held\_weights.sum()}.

\item[Q: Does \texttt{run\_backtest} produce plots of how contributions/weights evolve, or does that happen elsewhere?]
Elsewhere. \texttt{run\_backtest} is a pure computation engine -- it
returns data (\texttt{BacktestResult}), never a figure, by design
(\texttt{CLAUDE.md}'s separation of concerns). Plotting belongs in
\texttt{riskparity/analysis/plots.py} (Section~\ref{sec:metrics}), not
yet implemented, which will consume one or more
\texttt{BacktestResult}s and save figures to \texttt{results/figures/}
for inclusion here via the \texttt{IfFileExists} guards already wired
into Section~\ref{sec:results}.

\item[Q: Add tests for what was just implemented in \texttt{run\_backtest}.]
Added \texttt{tests/test\_backtest.py}, 11 tests, listed in
Section~\ref{sec:backtest-impl}'s ``Tests'' paragraph -- covering
\texttt{\_rebalance\_dates}'s resample/dropna behavior, the warm-up
skip logic, structural invariants on \texttt{weights\_history} and
\texttt{risk\_contributions\_history}, an explicit no-look-ahead check
via a spy \texttt{strategy\_fn}, and exact (not approximate) turnover
and transaction-cost accounting.

\item[Q: What are the next steps for the project?]
Cross-checked \texttt{CLAUDE.md}'s suggested build order against actual
module contents: steps 1--7 (volatility, vol-ERC, tail risk, Component
CVaR, CVaR-parity, data download/clean, backtest engine) were all
implemented and the full suite passed (50/50 tests at that point).
Remaining: step 8, \texttt{analysis/compare.py} + \texttt{plots.py}
(both still \texttt{raise NotImplementedError}); step 9, wiring
everything into \texttt{scripts/run\_pipeline.py}; step 10, filling in
the README's ``Key results'' section with real numbers. Agreed order:
\texttt{compare.py} first (its output is what \texttt{plots.py} and the
pipeline both need), then \texttt{plots.py}, then the pipeline wiring.

\item[Q: Start with \texttt{analysis/compare.py}.]
Implemented \texttt{performance\_table} and
\texttt{risk\_contribution\_stability}, both documented in full in
Section~\ref{sec:compare-impl}. \texttt{diversification\_ratio}, also
mentioned in the module's original TODO, was deliberately left
unimplemented and flagged rather than silently dropped or worked around
-- it needs per-asset volatilities that \texttt{BacktestResult} does
not currently expose (see that section's ``Deferred'' paragraph).
Along the way, the test suite for \texttt{performance\_table} caught a
real bug: a zero-volatility edge case produced an astronomically large
Sharpe ratio instead of the intended \texttt{NaN}, because float64
summation noise left the computed volatility at $\sim 10^{-19}$ instead
of exactly \texttt{0.0}, which passed a naive \texttt{> 0} guard. Fixed
with an explicit $10^{-12}$ epsilon threshold, matching the numerical
floor pattern already used in \texttt{optimization/erc.py}.

\item[Q: Now implement \texttt{analysis/plots.py}.]
Implemented the three functions already declared in the module
(\texttt{plot\_risk\_contributions}, \texttt{plot\_cumulative\_returns},
\texttt{plot\_contribution\_stability}), documented in full in
Section~\ref{sec:plots-impl}. The module docstring's fourth suggested
figure (a stress-window zoom) was deliberately left without its own
function, since it is exactly \texttt{plot\_cumulative\_returns} called
on pre-sliced data -- same reasoning as deferring
\texttt{diversification\_ratio} in \texttt{compare.py}: flag a gap
explicitly rather than add a near-duplicate function or silently drop
the requirement. Sanity-checked end-to-end (not just via
\texttt{tests/test\_plots.py}'s structural checks) by running
\texttt{solve\_erc} and \texttt{solve\_cvar\_parity} through
\texttt{run\_backtest} on a synthetic universe and visually inspecting
the three saved PNGs: results matched expectations, including
CVaR-parity's known non-convergence showing up as real, visible
contribution drift in \texttt{plot\_contribution\_stability} where
vol-ERC (on an i.i.d.\ synthetic universe) stayed essentially flat.

\end{description}

% ============================================================================
\section{Glossary of Finance Concepts}
\label{sec:glossary}
% ============================================================================

\begin{description}[leftmargin=3.5cm, style=nextline]
    \item[Risk parity / ERC] Portfolio construction approach that
    allocates so that each asset contributes equally to total portfolio
    risk, rather than allocating equal capital (equal-weight) or
    minimizing total risk regardless of concentration (min-variance).
    \item[Diversification ratio] Ratio of the weighted average of
    individual asset volatilities to portfolio volatility; equals~1 for a
    single asset and increases with genuine diversification benefit.
    \item[Value at Risk (VaR)] The loss level not expected to be exceeded
    with probability $\alpha$ over a given horizon; a quantile of the
    loss distribution.
    \item[Conditional VaR (CVaR) / Expected Shortfall] The expected loss
    \emph{given} that the loss exceeds VaR; averages the tail beyond the
    VaR quantile rather than reporting a single cutoff, and is a coherent
    (subadditive) risk measure, unlike VaR.
    \item[Tail risk] Risk of extreme, low-probability losses that are
    underestimated by measures (like volatility) that only use the first
    two moments of the return distribution.
    \item[Covariance shrinkage] Regularizing a sample covariance matrix
    estimate toward a structured target (e.g.\ constant correlation) to
    reduce estimation error, especially important when the number of
    assets is not small relative to the number of observations.
    \item[Turnover] Fraction of portfolio value traded at a rebalance;
    drives transaction costs.
    \item[Drawdown] Peak-to-trough decline in cumulative portfolio value;
    \emph{maximum drawdown} is the worst such decline over the backtest.
    \item[Sharpe ratio] Annualized excess return divided by annualized
    volatility; a risk-adjusted return measure that does \emph{not}
    distinguish upside from downside volatility (unlike CVaR-based
    measures).
    \item[Calmar ratio] Annualized return divided by maximum drawdown; a
    tail-oriented alternative to the Sharpe ratio.
\end{description}

% ============================================================================
\begin{thebibliography}{9}

\bibitem{maillard2010}
Maillard, S., Roncalli, T., \& Teïletche, J. (2010).
\emph{The Properties of Equally Weighted Risk Contribution Portfolios.}
Journal of Portfolio Management.

\bibitem{boudt2013}
Boudt, K., Peterson, B., \& Carl, P. (2013).
\emph{Component VaR and CVaR.}

\bibitem{rockafellar2000}
Rockafellar, R.~T., \& Uryasev, S. (2000).
\emph{Optimization of Conditional Value-at-Risk.}
Journal of Risk.

\bibitem{roncalli2013}
Roncalli, T. (2013).
\emph{Introduction to Risk Parity and Budgeting.}
Chapman \& Hall/CRC.

\end{thebibliography}

\end{document}
